
\section{Введение}

Память это одна из высших психических функций, позволяющая сохранять,
менять и бесконечно воспроизводить информацию. \cite{wikipedia_memory_ru}
Люди используют её постоянно, чтобы запоминать разговоры с людьми,
вспоминать о прошлом или же творить, основываясь на уже увиденном. Будь
то книги, фильмы, музыка или живопись, каждый раз добавляя свой личный
опыт, накопленный за всё время жизни.

Эта функция мозга притягивает многих исследователей.
Одним из них был Джон Джозеф Хопфилд (Рис. \ref{media:hopfield}), получивший в 2024 году
Нобелевскую Премию по физике за “Фундаментальные открытия и изобретения,
которые позволяют осуществлять машинное обучение с использованием
искусственных нейронных сетей.” \cite{wikipedia_hopfield_en}

\begin{wrapfigure}{r}{0.25\textwidth}
  \centering
  \includegraphics[width=0.25\textwidth]{image_hopfield.png}
  \caption{}
  \label{media:hopfield}
\end{wrapfigure}

В частности, в своей работе 1982 года ``Нейронные сети и физические
системы с возникающими коллективными вычислительными способностями'' он
ввёл понятие нейронной сети Хопфилда, использующей бинарные нейроны,
которая может быть ассоциативной памятью. \cite{hopfield1982neural}
То есть может сохранять и воспроизводить информацию.

Также в 1984 году появилась непрерывная классическая модель Хопфилда в
работе ``Нейроны с градуированным ответом обладают коллективными
вычислительными свойствами, подобными свойствам двухсостоятельных
нейронов.'', которая описывает возможность использования непрерывного
состояния и динамики. \cite{hopfield1984neurons}

Однако, в 2020 группой учёных была создана современная модель Хопфилда с
добавлением скрытого слоя, что позволило увеличить способность нейросети
до хранения экспоненциального количества паттернов. \cite{ramsauer2020hopfield}
Её часто упоминают при разговорах о механизмах внимания современных
трансформеров, которая обеспечивает логическую связность и целостность
результата генерации будь то текста или изображения.

\subsection{\texorpdfstring{\textbf{
      Постановка задачи}}
  {Постановка задачи}}
\textbf{Цель исследования:}
применимость скоростного градиента к обучению
классической непрерывной сети Хопфилда.

\textbf{Задачи исследования:}
\begin{itemize}
  \item
        Реализовать классическую и непрерывную модели Хопфилда;
  \item
        Провести обучение с помощью правила Хэбба, Сторки и скоростного
        градиента;
  \item
        Продемонстрировать результаты с разными функциями шума и
        параметром регуляризации;
  \item
        Применить непрерывную сеть Хопфилда с обучением скоростным градиентом
        к задаче LMAPF.
\end{itemize}